% jacobiano.tex
%
% Copyright (C) 2022--2026 José A. Navarro Ramón <janr.devel@gmail.com>
% Licencia del código GPLv2
% Licencia Creative Commons Recognition Non-Commercial Share-alike.
% (CC-BY-NC-SA)

\chapter{Jacobiano de una transformación}
\section{Cambio de variable en integración de una variable}
Normalmente se utiliza un cambio de variable (sustitución) cuando queremos simplificar
una integral.
Supongamos que tenemos la siguiente integral
\begin{equation}
  I= \int_a^b \!f(x)\, dx
\end{equation}
y, para resolverla, realizamos la siguiente sustitución $x = g(u)$.
Como consecuencia, obtenemos también la diferencial $dx=g'(u)\,du$.
\begin{equation}
  I = \int_c^d \!f(g(u))\,g'(u)\,du = \int_c^d\!(g\circ f)(u)\,g'(u)\,du
\end{equation}
donde $a=g(c)$ y $b=g(d)$. La transformación debe ser, al menos $C^1$, esto es, debe
existir la primera derivada y ser continua.

La transformación también se puede escribir como
\begin{align}
  x &= x(u)\\
  dx &= x'(u)\,du
\end{align}
y entonces, la integral quedaría
\begin{equation}
  I = \int_c^d\!f(g(u))\,\dfrac{dx}{du}\,du
\end{equation}
Vemos que para resolver la integral, tenemos que sustituir en $f(x)$ la $x$ por $u$, y
después sustituir $dx$ por $dx/du$ multiplicado por $du$
\begin{equation}
  dx = \dfrac{dx}{du}\,du
\end{equation}
Esto se puede interpretar como que, al cambiar de variable, de $x$ a $u$, mediante una
transformación $x=g(u)$ o $u=g^{-1}(x)$, la diferencial de longitud $dx$ se puede
sustituir por una derivada multiplicada por la diferencial de la nueva variable.
Aunque sea adelantar conceptos, y ser muy poco precisos, lamamos
\emph{jacobiano de la transformación} a lo que multiplica a la diferencial de la nueva
variable.
En este ejemplo de una variable, el jacobiano es una derivada (determinante $1\times 1$).
En integrales múltiples de orden $N$ el jacobiano será un determinante $N\times N$.





%%% Local Variables:
%%% mode: latex
%%% TeX-engine: luatex
%%% TeX-master: "../retazosmatematicas.tex"
%%% End:
